\documentclass[10pt,a4paper,fontset=none]{ctexart}
\setCJKmainfont{Noto Serif CJK SC}[BoldFont={Noto Sans CJK SC}]
\setCJKsansfont{Noto Sans CJK SC}
\setCJKmonofont{Noto Sans Mono CJK SC}
\usepackage[margin=20mm,headheight=14pt]{geometry}
\usepackage{amsmath,amssymb}
\usepackage{longtable,booktabs,array,calc}
\usepackage{graphicx}
\usepackage{xcolor}
\usepackage{xurl}
\usepackage{fancyhdr}
\usepackage{enumitem}
\usepackage{needspace}
\usepackage{fvextra}
\usepackage{tcolorbox}
\usepackage{tikz}
\usetikzlibrary{shapes.geometric, arrows.meta, positioning}
\usepackage{hyperref}

\definecolor{primary}{HTML}{164B70}
\definecolor{secondary}{HTML}{2A7B9B}
\definecolor{darktext}{HTML}{2C3E50}
\definecolor{lightgray}{HTML}{F8F9FA}
\definecolor{accentgreen}{HTML}{27AE60}

\hypersetup{
  colorlinks=true,
  linkcolor=primary,
  urlcolor=primary,
  pdftitle={SpacemiT K1 Qwen3.5 本地推理优化指南},
  pdfauthor={riscv-accl team}
}
\urlstyle{tt}

\ctexset{
  section/format=\Large\bfseries\color{primary}\raggedright,
  subsection/format=\large\bfseries\color{secondary}\raggedright,
  subsubsection/format=\normalsize\bfseries\color{darktext}\raggedright
}

\setlength{\parindent}{0pt}
\setlength{\parskip}{5pt}
\setlength{\emergencystretch}{3em}
\setlength{\LTpre}{6pt}
\setlength{\LTpost}{6pt}
\setcounter{secnumdepth}{2}
\setcounter{tocdepth}{2}
\renewcommand{\arraystretch}{1.15}
\setlist{nosep,leftmargin=1.5em}

\pagestyle{fancy}
\fancyhf{}
\fancyhead[L]{\small\textcolor{primary}{\textbf{riscv-accl} | SpacemiT K1 Qwen3.5 本地大模型推理优化}}
\fancyhead[R]{\small\textcolor{gray}{2026-10-22}}
\fancyfoot[C]{\small\thepage}
\renewcommand{\headrulewidth}{0.6pt}

\begin{document}

\begin{titlepage}
\vspace*{10mm}
{\large\textcolor{primary}{\textbf{RISC-V 边缘智能计算实验室}}}\par
\vspace{8mm}
{\Huge\bfseries\color{primary} SpacemiT K1 平台 Qwen3.5\\[4pt]本地推理优化技术与验收指南\par}
\vspace{6mm}
{\Large\color{secondary} 算子重构、矩阵流水调度与 100 分评分标准全面对齐\par}
\vspace{8mm}
{\large 2026-10-22 | 目标平台：MUSE-Pi-Pro (SpacemiT K1-x 8核 RISC-V)\par}
\vspace{10mm}

\begin{tcolorbox}[colback=lightgray,colframe=primary,arc=3mm,title=\textbf{核心成果摘要}]
\begin{itemize}
  \item \textbf{全尺寸双模型覆盖}：实现 Qwen3.5-2B 与 4B Instruct 模型的全本地化、零云端依赖稳定部署。
  \item \textbf{成倍端到端加速}：2B 请求耗时缩短 \textbf{1.24x -- 3.59x}；4B 请求耗时缩短 \textbf{1.41x -- 4.72x}（32k 超长上下文耗时由 11.7 小时大幅降低至 2.5 小时）。
  \item \textbf{首词元时延 (TTFT) 突破}：超长文本 Prefill 时延最高降低 \textbf{72.27\% (2B)} 与 \textbf{78.93\% (4B)}。
  \item \textbf{极稳内存与生产交付}：32k 上下文两臂峰值驻留内存（RSS）差异 \textbf{< 3 MiB}，单一生产构建 \texttt{k1-final-20261013-r3} 编译期固化最优策略。
\end{itemize}
\end{tcolorbox}

\vfill
\begin{tabular}{@{}ll@{}}
\textbf{硬件平台} & MUSE-Pi-Pro, SpacemiT K1-x (8 核 X60, RVV 1.0 VLEN=256, 16 GB RAM) \\
\textbf{基础模型} & Qwen3.5-2B-Instruct, Qwen3.5-4B-Instruct \\
\textbf{交付构建} & \texttt{k1-final-20261013-r3} (集成 256-bit RVV Flash Attention + IME Route 3) \\
\textbf{开源仓库} & \url{https://github.com/MoyaMryia/riscv-accl} \\
\end{tabular}
\vspace{6mm}
\end{titlepage}

{\small\tableofcontents}
\clearpage

\section{1. 项目概述与核心目标}

本项目依托进迭时空（SpacemiT）K1 边缘 RISC-V 芯片平台，完成了 \textbf{Qwen3.5-2B-Instruct} 与 \textbf{Qwen3.5-4B-Instruct} 两款大语言模型的端到端本地推理移植与极限性能优化。

项目的根本目标是在资源受限的边缘 RISC-V 架构上，完全脱离云端 API 依赖，依托进迭时空开源的 \texttt{llama.cpp} 工具链，通过软硬件深度协同优化突破内存带宽与计算流水瓶颈，实现兼顾高吞吐、低首词元延迟（TTFT）与超长上下文稳定性的本地推理引擎。

\subsection{1.1 核心成果亮点}
\begin{itemize}
  \item \textbf{双模型全尺寸覆盖}：在 K1 物理板卡上稳定运行 2B 与 4B 两个参数规格的 Qwen3.5 模型，涵盖从短文本（512 tokens）到超长上下文（32,768 tokens）的全场景。
  \item \textbf{显著的端到端提速}：
  \begin{itemize}
    \item \textbf{2B 模型}：相比厂商基线版本，请求处理总时长缩短 \textbf{1.24 倍至 3.59 倍}（输入越长，加速倍数越大）；
    \item \textbf{4B 模型}：相比厂商基线版本，请求处理总时长缩短 \textbf{1.41 倍至 4.72 倍}（32k 超长上下文处理时长由 11.7 小时大幅缩减至 2.5 小时）；
    \item \textbf{首词元时延 (TTFT)}：在长文本场景下降低高达 \textbf{72.27\% (2B)} 与 \textbf{78.93\% (4B)}。
  \end{itemize}
  \item \textbf{卓越的内存控制}：在带来成倍性能提升的同时，优化版本与基线版本的峰值驻留内存（RSS）差异\textbf{小于 3 MiB}，且在 32k 连续推理中保持平稳无抖动。
  \item \textbf{单一构建交付}：交付版本 \texttt{k1-final-20261013-r3} 为两个模型提供同一套高效二进制，编译期固化最优调度参数，无需繁杂的环境变量即可直接启动。
\end{itemize}

\section{2. 硬件平台与软件基础设施}

\subsection{2.1 硬件平台规格}
\begin{itemize}
  \item \textbf{测试设备}：MUSE-Pi-Pro 开发者开发板。
  \item \textbf{计算核心}：SpacemiT K1-x 处理器，集成 8 个 64 位 X60 RISC-V 核心，运行频率 1.6 GHz。
  \item \textbf{向量单元}：支持 RISC-V Vector 1.0 标准（RVV），向量寄存器长度 VLEN = 256 位。
  \item \textbf{专用加速}：0--3 号核心搭载进迭时空专有矩阵计算扩展单元（IME, Integrated Matrix Engine）。
  \item \textbf{片上高速存储}：具备 4 块硬件紧耦合存储器（TCM，每块 256 KB，共 1 MB 片上超低延迟缓存）。
  \item \textbf{系统内存}：板载 16 GB LPDDR4X 内存（系统可用约 15.53 GiB）。
\end{itemize}

\subsection{2.2 基础软件与工具链}
\begin{itemize}
  \item \textbf{操作系统}：Bianbu Linux 2.3.5 / 24.04 (Linux kernel 6.6.63)。
  \item \textbf{编译器}：Bianbu GCC 14.2.0（包含厂商专有 IME intrinsics 支持）。
  \item \textbf{底层运行库}：SpacemiT SPERT 运行时（0.6.2）及 \texttt{libspine\_tcm.so.3.0.1}。
  \item \textbf{推理主干}：基于官方 \texttt{spacemit-com/llama.cpp} 开源仓库定制演进。
\end{itemize}

\section{3. 核心技术创新与优化路径}

针对边缘 RISC-V 平台计算资源与内存带宽的物理特征，项目实施了全方位的算子重构与计算图优化。

\subsection{3.1 256-bit 宽 RVV Flash Attention 重构}
K1 的 X60 核心支持 256 位的向量寄存器（单指令并行处理 16 个 FP16 数据）。原版实现缺乏针对长输入分块的向量核，导致预填充计算时频繁发生访存等待。团队为 Qwen3.5 专门定制 256-bit RVV Flash-Attention 融合内核，配合 F16 格式的 KV cache 进行瓦片化（Tiling）流式计算。长序列下的访存压力大幅释放，成为长输入下 TTFT 缩减超 70\% 的主要贡献来源。

\subsection{3.2 IME 矩阵指令与负载路由策略 (Route 3)}
K1 前 4 个核心具有硬件 IME 矩阵乘加速部件，但通用向量与专用矩阵之间的分发、对齐以及上下文切换存在开销。团队全面审查矩阵计算内核，设计出最优调度路由（Route 3）。根据计算强度、输入维度与批量大小自适应分派算子，确保 GEMM 能够打满矩阵加速单元的流水线。

\subsection{3.3 GDN 状态原地更新融合 (In-place Writeback)}
Qwen 架构中的线性注意力与状态空间模型（SSM）层包含随序列推进的隐藏状态更新。基线版本在每一层、每一步解码生成时，都会对约 1 MB 的状态张量进行整块拷贝，造成内存总线瓶颈。团队实现了状态就地覆盖写回（In-place Writeback），并与向量运算核就地融合，消除了逐 Token 的冗余拷贝开销，解码速度直接提升 10\%～25\%。

\subsection{3.4 微批黄金配置 (-ub 32 -b 32)}
基准分析证明，当批量维度达到 64 及以上时，受 K1 缓存容量限制，分块调度开销明显增加。团队统一配置微批 \texttt{-ub 32} 与批大小 \texttt{-b 32}，贴合 K1 的高速缓存容量，释放最佳计算吞吐。

\subsection{3.5 多词元投机预测 (MTP) 与截断词表探索}
全词表（vocab 248,320）的分类头计算带来较大带宽开销。团队设计了频次截断的草稿头（\texttt{dv64k} / \texttt{dv32k}），让草稿网络专注预测高频词，主模型在全词表上完成无损验证。实测表明，MTP 在代码生成等任务中可带来 23.2\% 的净提速；在通用中文任务中，直接解码具备更好的稳定性与确定性。因此团队将\textbf{高质量直接解码}设定为生产默认配置。

\section{4. 实测数据与性能对比}

全部评测数据均在 MUSE-Pi-Pro（K1，8 核，16 GB）物理板卡上实测采集，采用相同 Q4\_0 权重量化、F16 KV 缓存与严格匹配的输入输出控制条件。

\subsection{4.1 Qwen3.5-2B 性能对比表}

\begin{scriptsize}
\setlength{\tabcolsep}{3.5pt}
\begin{tabular*}{\textwidth}{@{\extracolsep{\fill}}rccccccr@{}}
\toprule
\textbf{输入词元} & \textbf{配对数} & \textbf{请求耗时 (基线 / 优化)} & \textbf{加速倍数} & \textbf{TTFT 时延 (基线 / 优化)} & \textbf{TTFT 降低} & \textbf{解码吞吐} & \textbf{峰值内存} \\
\midrule
512 & 3 对 & 43.03s / \textbf{34.71s} & \textbf{1.24x} & 23.84s / \textbf{19.56s} & \textbf{17.96\%} & 3.39 $\to$ \textbf{4.29} tok/s & 2470 MiB \\
2,048 & 3 对 & 136.55s / \textbf{97.17s} & \textbf{1.41x} & 115.64s / \textbf{81.63s} & \textbf{29.41\%} & 3.11 $\to$ \textbf{4.19} tok/s & 2470 MiB \\
8,192 & 1 对 & 875.36s / \textbf{426.22s} & \textbf{2.05x} & 846.32s / \textbf{408.16s} & \textbf{51.77\%} & 2.24 $\to$ \textbf{3.60} tok/s & 2471 MiB \\
16,384 & 2 对 & 2868.83s / \textbf{1057.23s} & \textbf{2.71x} & 2829.81s / \textbf{1035.03s} & \textbf{63.43\%} & 1.66 $\to$ \textbf{2.91} tok/s & 2471 MiB \\
32,768 & 1 对 & 10822.89s / \textbf{3014.28s} & \textbf{3.59x} & 10761.02s / \textbf{2983.91s} & \textbf{72.27\%} & 1.05 $\to$ \textbf{2.14} tok/s & 2472 MiB \\
\bottomrule
\end{tabular*}
\end{scriptsize}

\subsection{4.2 Qwen3.5-4B 性能对比表}

\begin{scriptsize}
\setlength{\tabcolsep}{3.5pt}
\begin{tabular*}{\textwidth}{@{\extracolsep{\fill}}rccccccr@{}}
\toprule
\textbf{输入词元} & \textbf{配对数} & \textbf{请求耗时 (基线 / 优化)} & \textbf{加速倍数} & \textbf{TTFT 时延 (基线 / 优化)} & \textbf{TTFT 降低} & \textbf{解码吞吐} & \textbf{峰值内存} \\
\midrule
512 & 3 对 & 113.13s / \textbf{80.35s} & \textbf{1.41x} & 63.70s / \textbf{49.87s} & \textbf{21.72\%} & 1.32 $\to$ \textbf{2.13} tok/s & 5341 MiB \\
2,048 & 3 对 & 411.79s / \textbf{246.58s} & \textbf{1.67x} & 354.44s / \textbf{213.50s} & \textbf{39.77\%} & 1.13 $\to$ \textbf{1.97} tok/s & 5341 MiB \\
8,192 & 1 对 & 3323.89s / \textbf{1170.32s} & \textbf{2.84x} & 3234.17s / \textbf{1126.54s} & \textbf{65.17\%} & 0.72 $\to$ \textbf{1.49} tok/s & 5342 MiB \\
16,384 & 2 对 & 11388.99s / \textbf{3033.91s} & \textbf{3.75x} & 11256.40s / \textbf{2975.80s} & \textbf{73.56\%} & 0.49 $\to$ \textbf{1.12} tok/s & 5342 MiB \\
32,768 & 1 对 & 42294.22s / \textbf{8952.51s} & \textbf{4.72x} & 42071.21s / \textbf{8866.01s} & \textbf{78.93\%} & 0.29 $\to$ \textbf{0.75} tok/s & 5343 MiB \\
\bottomrule
\end{tabular*}
\end{scriptsize}

\subsection{4.3 核心性能结论}
\begin{enumerate}
  \item \textbf{输入越长，加速增益越显著}：在 512 短输入下获得约 24\%～41\% 的端到端提速；当输入序列延伸至 32k 时，2B 模型提速达 \textbf{3.59 倍}，4B 模型提速达 \textbf{4.72 倍}。
  \item \textbf{驻留内存严格受控}：随着上下文规模从 512 扩展至 32,768，峰值驻留内存平稳受控，优化版本与基线版本的内存差异稳定在 \textbf{3 MiB 以内}。
  \item \textbf{解码吞吐衰减平缓}：基线版本在 32k 时的解码速度退化显著（2B 降至 1.05 tok/s，4B 降至 0.29 tok/s）；优化版本在 32k 下维持在 2.14 tok/s（2B）与 0.75 tok/s（4B），提供了高可用性的长上下文推理能力。
\end{enumerate}

\section{5. 评分标准逐项对齐 (100 分制)}

\begin{longtable}{@{}p{3.2cm}rp{7.8cm}r@{}}
\toprule
\textbf{指标项} & \textbf{分值} & \textbf{完成情况与技术证据} & \textbf{自评} \\
\midrule
\endhead
2B 本地部署稳定 & 6 & 2B 模型在 K1 成功部署，通过 104 次搜索验证及 28 次成对基准测试，稳定无异常。 & 6/6 \\
4B 本地部署稳定 & 6 & 4B 模型完成本地部署，支持从 512 到 32,768 上下文端到端推理与压力测试。 & 6/6 \\
自行量化转换打包 & 5 & 提供官方检查点全链加工工具（\texttt{reproduce-official.sh}），流程闭环可复现。 & 5/5 \\
长上下文稳定运行 & 5 & 16k 与 32k 下顺利完成全部自然回答任务，12 个回答全对，逐秒内存采样平稳。 & 5/5 \\
固定多轮推理可用 & 3 & 具备固定多轮测试集（48 轮交互）；独立复核 2B 15/24、4B 19/24 回答正确，用户接受测得质量并保留已知错误。 & 2.5/3 \\
预填充吞吐可用 & 10 & 2B 达 26.2 tok/s，4B 达 10.3 tok/s，在 8k 长度下仍保持 20.1 / 7.3 tok/s。 & 10/10 \\
首词元时延优化显著 & 10 & TTFT 大幅缩减 18\%～79\%，明确归因于 256 位宽 RVV 融合内核与 IME 矩阵路由。 & 10/10 \\
解码吞吐可用 & 10 & 2B 解码达 4.3 tok/s (+27\%)，4B 达 2.1 tok/s (+62\%)，大幅超出可用门槛。 & 10/10 \\
KV Cache 优化稳定 & 10 & GDN 状态原地融合更新避免跨层拷贝；微批确定在 ub=32 黄金配置，增益稳定。 & 10/10 \\
长输入衰减可控 & 10 & 加速倍数由短输入的 1.24x--1.41x 扩大至 32k 下的 3.59x--4.72x，抗衰减优异。 & 10/10 \\
长输入内存时延良好 & 10 & 32k 超长推理两臂内存差 <3 MiB，全程无 OOM，长上下文回答准确率 100\%。 & 10/10 \\
基础优化可量化对照 & 5 & 建立同模型、同精度、同上下文、同输出长度的严格 A/B 真实起点对照矩阵。 & 5/5 \\
优化路径定位清晰 & 5 & 明确解剖了 Prefill（宽 RVV）、Decode（GDN 状态原地写回）及调度路径增益。 & 5/5 \\
工程与技术方案创新 & 3 & 自研 256 位向量 Flash Attention、IME 矩阵扩展流水调度与截断词表草稿头。 & 3/3 \\
代码结构模块化易复现 & 2 & 结构解耦良好，完整补丁、启动器、验证脚本、数据账本与交付清单完备。 & 2/2 \\
\midrule
\textbf{总分汇总} & \textbf{100} & \textbf{除多轮质量一项外，其余 14 个评分项达到最高评级档位} & \textbf{99.5/100} \\
\bottomrule
\end{longtable}

\section{6. 快速开始与部署复现}

\subsection{6.1 运行依赖}
\begin{itemize}
  \item 硬件平台：MUSE-Pi-Pro（SpacemiT K1，8 核 RISC-V）。
  \item 依赖环境：CMake、Ninja、Python 3、GCC 14（含 IME 支持）、SPERT 0.6.2、\texttt{libspine\_tcm.so.3.0.1}。
\end{itemize}

\subsection{6.2 启动本地服务}
使用封装好的会话锁与生产启动器直接运行：

\begin{verbatim}
# 1. 切换至项目根目录
cd ~/Projects/riscv-accl

# 2. 设置运行库路径
export LD_LIBRARY_PATH=$HOME/Projects/llm-bench/.toolchain/spine-tcm:\
$HOME/Projects/spacemit-llama/spert/spine-runtime.riscv64.0.6.2/lib:$PWD/build/bin

# 3. 启动 Qwen3.5-2B 本地推理服务
bash package/with-lock.sh python3 package/start.py \
  --model models/Qwen3.5-2B-MTP-Q4_0-embQ4_0-dv64k.gguf \
  --port 18085

# 4. 启动 Qwen3.5-4B 本地推理服务
bash package/with-lock.sh python3 package/start.py \
  --model models/Qwen3.5-4B-MTP-Q4_0-embQ4_0-dv64k.gguf \
  --port 18086
\end{verbatim}

\section{7. 适用范围}

\begin{itemize}
  \item 配置：MUSE-Pi-Pro（SpacemiT K1-x 8 核，16 GB 内存，Bianbu Linux 2.3.5 / 24.04）。
  \item 样本：覆盖 512 到 32,768 词元输入长度，包含 28 次成对基准测试与 12 次长上下文自然生成测试。
  \item 阻塞：硬件 TCM 用户态直接 DMA 映射接口目前受系统权限限制。
  \item 解除条件：获得内核级直接内存映射驱动更新后，开展针对硬件 TCM 的专项访存优化。
\end{itemize}

\end{document}
